\documentclass[10pt, letterpaper]{article}

% Packages:
\usepackage[
    ignoreheadfoot,
    top=2cm,
    bottom=2cm,
    left=2cm,
    right=2cm,
    footskip=1.0cm,
]{geometry}
\usepackage{titlesec}
\usepackage{tabularx}
\usepackage{array}
\usepackage[dvipsnames]{xcolor}
\definecolor{primaryColor}{RGB}{0, 79, 144}
\usepackage{enumitem}
% \usepackage{fontawesome5} % not used; commented out for portability
\usepackage{amsmath}
\usepackage[
    pdftitle={JaeHyun Ha's CV},
    pdfauthor={JaeHyun Ha},
    pdfcreator={LaTeX with RenderCV},
    colorlinks=true,
    urlcolor=primaryColor
]{hyperref}
\usepackage[pscoord]{eso-pic}
\usepackage{calc}
\usepackage{bookmark}
\usepackage{lastpage}
\usepackage{changepage}
\usepackage{paracol}
\usepackage{ifthen}
\usepackage{needspace}
\usepackage{iftex}
\usepackage{tabularx}
\usepackage{url}
\usepackage{array}

\ifPDFTeX
    \input{glyphtounicode}
    \pdfgentounicode=1
    \usepackage[utf8]{inputenc}
    \usepackage{lmodern}
\fi

\pagestyle{empty}
\setcounter{secnumdepth}{0}
\setlength{\parindent}{0pt}
\setlength{\columnsep}{0cm}

% Section formatting
\titleformat{\section}{\needspace{4\baselineskip}\bfseries\large}{}{0pt}{}[\vspace{1pt}\titlerule]
\titlespacing{\section}{-1pt}{0.3cm}{0.2cm}

% Environments
\renewcommand\labelitemi{$\circ$}
\newenvironment{highlights}{
    \begin{itemize}[
        topsep=0.10cm,
        parsep=0.10cm,
        partopsep=0pt,
        itemsep=0pt,
        leftmargin=0.4cm+10pt
    ]
}{
    \end{itemize}
}
\newenvironment{onecolentry}{
    \begin{adjustwidth}{0.2cm}{0.2cm}
}{
    \end{adjustwidth}
}
\newenvironment{twocolentry}[2][]{
    \onecolentry
    \def\secondColumn{#2}
    \setcolumnwidth{\fill, 4.5cm}
    \begin{paracol}{2}
}{
    \switchcolumn \raggedleft \secondColumn
    \end{paracol}
    \endonecolentry
}

\begin{document}

\begin{center}
    {\LARGE \textbf{Jaehyun Ha}} \\[5pt]
\end{center}

\begin{tabularx}{\textwidth}{X r}
Department of Artificial Intelligence, POSTECH, South Korea &
\href{https://dslab.postech.ac.kr/}{Data Systems Lab @ POSTECH}
\\
Homepage: \href{https://githajae.github.io}{githajae.github.io} &
Email: \href{mailto:jhha@dblab.postech.ac.kr}{jhha@dblab.postech.ac.kr}
\\
Advisor: \href{https://wscrony.github.io/}{Wook-Shin Han} \\
\end{tabularx}

\vspace{0.4cm}

\section*{Research Interests}

\vspace{1mm}
\noindent
\textbf{Building AI-native database systems for unified analytics and inference on multi-modal data}

\begin{itemize}
    \itemsep0.5em
    \item \textbf{Graph Databases:} Architecting a foundational query framework based on a unified graph representation of multi-modal data

    \item \textbf{Semantic Operator:} Embedding Large Language Models (LLMs) into the database to enable context-aware querying of unstructured data

    \item \textbf{Query Optimization:} Designing proxy models and novel cardinality estimation methods to optimize the execution of computationally expensive semantic operators
\end{itemize}

\section{Education}
\begin{twocolentry}{Mar. 2022 -- Present \;\; (Expected: Feb. 2028)}
\textbf{POSTECH}, Graduate School of Artificial Intelligence\\
M.S./Ph.D. Integrated Program\\
\textit{Data Systems Lab} \\
Advisor: Prof. Wook-Shin Han\\
\end{twocolentry}

\vspace{0.2cm}
\begin{twocolentry}{Mar. 2018 -- Mar. 2022}
\textbf{POSTECH}, Dept. of Computer Science and Engineering\\
B.S. in Computer Science and Engineering\\
GPA: 4.06 / 4.3\\
\textit{Graduated with the \textbf{highest} GPA in CS Dept.}
\end{twocolentry}

\vspace{0.2cm}
\begin{twocolentry}{Mar. 2016 -- Mar. 2018}
\textbf{Ulsan Science High School}
\end{twocolentry}


\section{Professional Experiences}

\begin{twocolentry}{Jun 2026 -- Sep 2026}
\textbf{Google Systems Research Group}, Sunnyvale, CA, USA
\end{twocolentry}
\begin{onecolentry}
\textit{Visiting Student Researcher}
\begin{highlights}
    \item Researching hardware-software co-optimization to accelerate semantic queries
    \item Hosted by Yeounoh Chung
\end{highlights}
\end{onecolentry}

\vspace{2mm}

\begin{twocolentry}{Jun 2025 -- Sep 2025}
\textbf{University of Illinois Urbana-Champaign (UIUC)}, IL, USA
\end{twocolentry}
\begin{onecolentry}
\textit{Visiting Scholar}
\begin{highlights}
    \item Conducted research on semantic operator optimization, resulting in a paper accepted at SIGMOD 2027
    \item Joint research project under the supervision of Prof. Yongjoo Park
\end{highlights}
\end{onecolentry}

\vspace{2mm}


\begin{twocolentry}{Jul 2020 -- Aug 2020}
\textbf{Digital Platform team}, SK hynix Inc., Korea
\end{twocolentry}
\begin{onecolentry}
\textit{Software Engineer Intern}
\begin{highlights}
    \item Developed real-time video comment overlay system for SK hynix's internal video streaming systems
    \item Won silver prize (2nd) in internship contest
\end{highlights}
\end{onecolentry}

\section*{Publications}

\newcommand{\myname}[1]{\textbf{#1}}

\begin{enumerate}[leftmargin=0.6cm, itemsep=1.5ex]
    \item \myname{Jaehyun Ha}, Yongjoo Park, Wook-Shin Han.
    \textit{CADENZA: Compiling Natural-Language Intent into Task-Specific Operator DAGs for Semantic Query Processing.} SIGMOD 2027.

    \item \myname{Jaehyun Ha}, Yongjoo Park, Wook-Shin Han.
    \textit{CADENZA in Action: Breaking the Monolith with Intent-Dependent Plan Spaces for Semantic Queries.} VLDB 2026 (Demo).

    \item Taesung Lee, \myname{Jaehyun Ha}, Byungchul Tak, Wook-Shin Han.
    \textit{TurboLynx: Schemaless Graph Engine Strikes Back for General-Purpose Analytics.} VLDB 2026.

    \item Taesung Lee, \myname{Jaehyun Ha}, Byungchul Tak, Wook-Shin Han.
    \textit{TurboLynx in Action: A Schemaless Graph Engine for General-Purpose Analytics.} VLDB 2026 (Demo).

    \item Wonseok Lee\textsuperscript{*}, \myname{Jaehyun Ha}\textsuperscript{*}, Wook-Shin Han, Changgyoo Park, Myunggon Park, Juhyeng Han, Juchang Lee. \textnormal{(\textsuperscript{*}equal contribution)}
    \textit{DoppelGanger++: Towards Fast Dependency Graph Generation for Database Replay.} SIGMOD 2024.

    \item Wonseok Lee, \myname{Jaehyun Ha}, Wook-Shin Han, Changgyoo Park, Myunggon Park, Juhyeng Han.
    \textit{DoppelGanger++ in Action: A Database Replay System with Fast Dependency Graph Generation.} VLDB 2024 (Demo).

    \item Kyoungmin Kim, \myname{Jaehyun Ha}, George Fletcher, Wook-Shin Han.
    \textit{Guaranteeing the O(AGM/OUT) runtime for uniform sampling and size estimation over joins.} PODS 2023.
\end{enumerate}

\section*{Invited Talks}

\begin{onecolentry}
\begin{itemize}[leftmargin=0.4cm]
    \item
    \begin{tabularx}{0.97\textwidth}{X r r}
        Optimizing Semantic Queries: Current Work and Open Questions & Cornell University & Sep 2026
    \end{tabularx}

    \item
    \begin{tabularx}{0.97\textwidth}{X r r}
        CADENZA: Intent-Dependent Planning for Semantic Queries & Google Systems Research Group & Jan 2026
    \end{tabularx}
\end{itemize}
\end{onecolentry}

\section*{Projects}

\begin{twocolentry}{Jun 2025 -- Sep 2025}
\textbf{Semantic Operator Optimization}
\end{twocolentry}
\begin{onecolentry}
\textit{Visiting Scholar Research at UIUC with Prof. Yongjoo Park}
\begin{highlights}
    \item \textbf{Problem:} Semantic operators are prohibitively expensive in latency and dollar cost. Existing optimization methods (i.e., proxy models) present a poor trade-off, forcing a choice between unacceptably low accuracy or still-significant latency and cost, which limits practical adoption
    \item \textbf{Contribution:} Developed novel methodologies to discover and build high-performance proxies that significantly improve latency and dollar cost while preserving accuracy
    \item \textbf{Outcome:} Accepted at SIGMOD 2027 and VLDB 2026 (Demo)
\end{highlights}
\end{onecolentry}

\vspace{2mm} % Spacing between projects

\begin{twocolentry}{Jan 2023 -- Jun 2025}
\textbf{High-Performance Schemaless Graph Database System}
\end{twocolentry}
\begin{onecolentry}
\begin{highlights}
    \item \textbf{Problem:} The conversion of unstructured data into knowledge graphs produces schemaless graphs, where nodes and edges have their own different schemas. However, existing database systems have significant limitations in performing high-performance analytics on such data
    \item \textbf{Contribution:} Designed and implemented a full-stack system with a specialized storage, optimizer, and query engine tailored for schemaless graph analytics
    \item \textbf{Role:} Core developer for a large-scale system (246K+ LOC); authored approx 50\% of commits
    \item \textbf{GitHub:} \url{https://github.com/postech-dblab-iitp/turbograph-v3}
    \item \textbf{Outcome:} Accepted at VLDB 2026 and VLDB 2026 (Demo)
\end{highlights}
\end{onecolentry}

\vspace{2mm}

\begin{twocolentry}{Dec 2022 -- Dec 2023}
\textbf{High-Speed Dependency Graph Generation for Database Replay}
\end{twocolentry}
\begin{onecolentry}
\textit{Industry Collaboration with SAP Labs Korea}
\begin{highlights}
    \item \textbf{Problem:} Database replay systems are innovative tools for capturing and replaying real-world workloads for testing purposes. However, their real-world efficiency is severely bottlenecked by the slow process of dependency graph generation
    \item \textbf{Contribution:} Proposed an efficient algorithm to accelerate dependency graph generation
    \item \textbf{Impact:} Implemented on the commercial SAP HANA system, demonstrating real-world utility
    \item \textbf{Outcome:} Accepted at SIGMOD 2024 and VLDB 2024 (Demo)
\end{highlights}
\end{onecolentry}

\vspace{2mm}

\begin{twocolentry}{Mar 2022 -- Nov 2022}
\textbf{Theoretically Optimal Join Cardinality Estimation Algorithm}
\end{twocolentry}
\begin{onecolentry}
\begin{highlights}
    \item \textbf{Problem:} Existing sampling-based algorithms for join cardinality estimation lacked a provable optimal O(AGM/OUT) runtime bound
    \item \textbf{Contribution:} Developed the first algorithm to achieve this theoretical optimal bound
    \item \textbf{Outcome:} Accepted at PODS 2023
\end{highlights}
\end{onecolentry}


\section*{Technical Skills}
\vspace{1mm}

\begin{onecolentry}
\renewcommand{\arraystretch}{1.18}
\setlength{\tabcolsep}{6pt}
\begin{tabularx}{\textwidth}{
    >{\raggedright\arraybackslash\bfseries}p{4.6cm}
    >{\raggedright\arraybackslash}X
}
Programming Languages & C/C++ (Advanced), Python (Advanced), SQL (Advanced), CUDA (Intermediate), Java (Intermediate), JavaScript (Intermediate) \\
Databases \& Data & PostgreSQL, OracleDB, DuckDB, Kuzu, Neo4j, pandas \\
Systems \& Build Tools & Linux, Git, Docker, Kubernetes, CMake, Bazel, uv \\
ML/LLM & PyTorch, Hugging Face, OpenAI API, Azure OpenAI API \\
\end{tabularx}
\end{onecolentry}


\section{Honors and Fellowships}
\begin{onecolentry}
\begin{itemize}[leftmargin=0.4cm]

    \item
    \begin{tabularx}{0.97\textwidth}{X r}
        POSTECHIAN Fellowship & 2022,2023
    \end{tabularx}

    \item
    \begin{tabularx}{0.97\textwidth}{X r}
        Samsung Dream Scholarship Foundation Scholar & 2019-2021
    \end{tabularx}

\end{itemize}
\end{onecolentry}

\section{Teaching Experiences}
\begin{onecolentry}
\begin{itemize}[leftmargin=0.4cm]

    \item
    \begin{tabularx}{0.97\textwidth}{X r}
        \textbf{Teaching Assistant}, Agentic RAG, Samsung Electronics & Spring 2026
    \end{tabularx}

    \item
    \begin{tabularx}{0.97\textwidth}{X r}
        \textbf{Teaching Assistant}, AIGS540: Big Data Processing, POSTECH & Spring 2024
    \end{tabularx}

    \item
    \begin{tabularx}{0.97\textwidth}{X r}
        \textbf{Teaching Assistant}, CSED421: Database Systems, POSTECH & Spring 2022
    \end{tabularx}

\end{itemize}
\end{onecolentry}

\section{References}

\begin{onecolentry}
    \begin{tabularx}{\textwidth}{@{} X r @{}}
        \textbf{Wook-Shin Han} & \href{mailto:wshan@dblab.postech.ac.kr}{wshan@dblab.postech.ac.kr} \\[2pt] % 2pt 정도 띄워줌
        \multicolumn{2}{@{} l @{}}{POSTECH Distinguished Professor, Pohang University of Science and Technology (POSTECH), Korea}
    \end{tabularx}
\end{onecolentry}

\vspace{0.2cm}

\begin{onecolentry}
    \begin{tabularx}{\textwidth}{@{} X r @{}}
        \textbf{Yongjoo Park} & \href{mailto:yongjoo@g.illinois.edu}{yongjoo@g.illinois.edu} \\[2pt]
        \multicolumn{2}{@{} l @{}}{Assistant Professor, Department of Computer Science, University of Illinois Urbana-Champaign (UIUC), USA}
    \end{tabularx}
\end{onecolentry}

\end{document}
